\section{Knowledge Retrieval and Purchase-Intent Inference}
\label{sec:inference}

\subsection{Inference inputs}

Service extracts structured items and attributes from the shared order and
assembles the context needed for intent judgment. This context may include
purchase-history signals, available calendar tags or temporal context, and
retrieved purchase-intent evidence. Execution-time device state, the device
catalog, and own-home execution memory are handled at their respective later
boundaries unless a further specification explicitly admits them as intent
evidence. \sourcecite{draft}

\subsection{Knowledge Base--Service contract}

The Knowledge Base stores generated purchase-intent cases and permitted
actual cases. Service queries for similar cases and receives a distribution
and representative evidence for runtime judgment and explanation. The
decision-time retrieval result preserves the following meanings.

\begin{itemize}
  \item The retrieval population used for the judgment.
  \item Intent and label identity associated with the returned evidence.
  \item Representative cases supporting the judgment.
  \item Provenance sufficient to reproduce the decision-time explanation.
\end{itemize}

These are shared information requirements. Exact serialized case fields,
receiver interpretation, and schema details remain under discussion.
\sourcecite{draft}

\open{Synthetic generation sources differ between joint full-case generation
and intent-seeded assembly. The shared draft preserves this conflict for a
team decision.}

\open{The mapping of generation-time confidence and ambiguity metadata into
the runtime case schema, including how Service receives and interprets that
metadata, remains unresolved.}

\subsection{Runtime judgment and task handoff}

Service performs one runtime purchase-intent judgment using the structured
order/context and retrieved evidence. Current semantics use structural and
format validation around that judgment. The resulting intent or no-intent
outcome is handed to the preparation stage. The baseline permits ordinary
shopping and previously rejected household situations to suppress preparation.
\sourcecite{draft}\sourcecite{baseline}

\open{Adoption of a runtime evidence or confidence threshold as an execution
gate requires a team decision. Offline dataset-quality thresholds and
runtime action policy have separate meanings. Dataset review criteria alone
do not establish runtime authorization.}

Runtime explanations use the evidence saved at decision time. Household
observation and user reactions later inform memory admission and future
preparation. The exact admission and observation-result semantics are covered
by the remaining cross-boundary decisions. \sourcecite{draft}
